\documentclass[11pt, a4paper]{article}
\usepackage{fontspec}\setmainfont{Helvetica Neue}\setmonofont{Menlo}
\usepackage{unicode-math}\setmathfont{STIX Two Math}
\usepackage[margin=2cm, top=2cm, bottom=2.5cm]{geometry}
\usepackage{microtype}\usepackage{parskip}
\setlength{\parskip}{0.6em}\setlength{\parindent}{0pt}
\usepackage[colorlinks=true, linkcolor=NavyBlue, urlcolor=NavyBlue, citecolor=NavyBlue]{hyperref}
\usepackage{xcolor}\usepackage[dvipsnames]{xcolor}
\usepackage{amsmath}\usepackage{booktabs}\usepackage{array}\usepackage{tabularx}
\usepackage{graphicx}
\graphicspath{{figures/week29/}}
\newcommand{\thead}[1]{\textbf{#1}}
\usepackage{enumitem}
\setlist[itemize]{leftmargin=1.5em, itemsep=0.2em, topsep=0.3em}
\setlist[enumerate]{leftmargin=1.5em, itemsep=0.2em, topsep=0.3em}
\usepackage[most]{tcolorbox}\tcbuselibrary{breakable, skins, theorems}
\definecolor{defBg}{HTML}{EAF2FB}\definecolor{defBd}{HTML}{1F4E79}
\definecolor{exBg}{HTML}{F4F4F4}\definecolor{exBd}{HTML}{555555}
\definecolor{tipBg}{HTML}{E8F5E9}\definecolor{tipBd}{HTML}{2E7D32}
\definecolor{trapBg}{HTML}{FCE4EC}\definecolor{trapBd}{HTML}{AD1457}
\definecolor{pitBg}{HTML}{FFF8E1}\definecolor{pitBd}{HTML}{B27600}
\definecolor{noteBg}{HTML}{F3E5F5}\definecolor{noteBd}{HTML}{6A1B9A}
\definecolor{tldrBg}{HTML}{E1F5FE}\definecolor{tldrBd}{HTML}{0277BD}
\newtcolorbox{defbox}[1][]{enhanced, breakable, colback=defBg, colframe=defBd, arc=2pt, boxrule=0.6pt, left=8pt, right=8pt, top=6pt, bottom=6pt, fonttitle=\bfseries, title={Definition: #1}}
\newtcolorbox{exbox}[1][]{enhanced, breakable, colback=exBg, colframe=exBd, arc=2pt, boxrule=0.6pt, left=8pt, right=8pt, top=6pt, bottom=6pt, fonttitle=\bfseries, title={Worked example: #1}}
\newtcolorbox{exambox}[1][]{enhanced, breakable, colback=tipBg, colframe=tipBd, arc=2pt, boxrule=0.6pt, left=8pt, right=8pt, top=6pt, bottom=6pt, fonttitle=\bfseries, title={Exam-mode answer #1}}
\newtcolorbox{trapbox}[1][]{enhanced, breakable, colback=trapBg, colframe=trapBd, arc=2pt, boxrule=0.6pt, left=8pt, right=8pt, top=6pt, bottom=6pt, fonttitle=\bfseries, title={MCQ trap #1}}
\newtcolorbox{pitfallbox}[1][]{enhanced, breakable, colback=pitBg, colframe=pitBd, arc=2pt, boxrule=0.6pt, left=8pt, right=8pt, top=6pt, bottom=6pt, fonttitle=\bfseries, title={Pitfall #1}}
\newtcolorbox{notebox}[1][]{enhanced, breakable, colback=noteBg, colframe=noteBd, arc=2pt, boxrule=0.6pt, left=8pt, right=8pt, top=6pt, bottom=6pt, fonttitle=\bfseries, title={Lecturer aside #1}}
\newtcolorbox{tldrbox}[1][]{enhanced, breakable, colback=tldrBg, colframe=tldrBd, arc=2pt, boxrule=0.6pt, left=8pt, right=8pt, top=6pt, bottom=6pt, fonttitle=\bfseries, title={TL;DR #1}}
\usepackage{titlesec}
\titleformat{\section}[block]{\Large\bfseries\color{defBd}}{\thesection.}{0.6em}{}
\titleformat{\subsection}[block]{\large\bfseries\color{defBd!80!black}}{\thesubsection}{0.5em}{}
\titleformat{\subsubsection}[block]{\normalsize\bfseries}{}{0pt}{}
\titlespacing*{\section}{0pt}{1.6em}{0.6em}\titlespacing*{\subsection}{0pt}{1.2em}{0.4em}
\usepackage{fancyhdr}\pagestyle{fancy}\fancyhf{}
\fancyhead[L]{\small CM52054 — Week 29}\fancyhead[R]{\small Kernel Methods}
\fancyfoot[C]{\small\thepage}\renewcommand{\headrulewidth}{0.3pt}
\newcommand{\examflag}{\textcolor{tipBd}{\textbf{[EXAM]}}\,}
\newcommand{\trapflag}{\textcolor{trapBd}{\textbf{[TRAP]}}\,}
\newcommand{\doflag}{\textcolor{defBd}{\textbf{[DO]}}\,}
\title{\vspace{-1cm}{\Huge\bfseries Week 29 — Kernel Methods in Machine Learning} \\[0.4em]{\large Formal kernel definition; building new kernels from old; polynomial, exponential, Gaussian; cosine kernel; comparison to deep learning}}
\author{\large CM52054 Foundational Machine Learning \\[0.2em]\normalsize University of Bath \,$\cdot$\, Dr Rohit Babbar}
\date{Revision notes}
\begin{document}\maketitle\thispagestyle{empty}\vspace{1em}

\begin{tldrbox}
\begin{itemize}[leftmargin=*]
  \item Wk29 makes the kernel notion from Wk26 rigorous and develops the \textbf{construction calculus}: which functions are kernels, and how to build new kernels from old ones.
  \item \textbf{Formal definition} \doflag{}: $k : \mathcal{X}\times\mathcal{X} \to \mathbb{R}$ is a kernel iff there exists a feature space $\mathcal{H}$ (a Hilbert space) and a feature map $\phi : \mathcal{X} \to \mathcal{H}$ such that $k(\mathbf{x}, \mathbf{x}') = \langle\phi(\mathbf{x}), \phi(\mathbf{x}')\rangle_{\mathcal{H}}$.
  \item \textbf{$\mathcal{X}$ doesn't need to be a vector space} — can be documents, graphs, strings, time series. Kernels generalise inner products to arbitrary input domains.
  \item \textbf{Feature map is non-unique.} The linear kernel $k(x, x') = \langle x, x'\rangle$ admits $\phi_1(x) = x$ in $\mathcal{H}_1 = \mathbb{R}$ and $\phi_2(x) = (x, x)/\sqrt 2$ in $\mathcal{H}_2 = \mathbb{R}^2$ — both give the same kernel.
  \item \textbf{Past paper Q2(6)}: Kernel methods can handle (a) high-dim data, (b) non-linear relationships, (c) neither, (d) both. \textbf{Answer: (d)} — both, simultaneously.
  \item \textbf{Construction rules} \doflag{} (memorise these — used in all kernel proofs):
  \begin{itemize}
    \item \textbf{Positive scalar multiple}: $\alpha k$ is a kernel for $\alpha > 0$.
    \item \textbf{Conic sum}: $\sum_j \alpha_j k_j$ is a kernel for $\alpha_j > 0$.
    \item \textbf{Product}: $k_1 \cdot k_2$ is a kernel.
    \item \textbf{Mapping wrap}: $\hat k(\mathbf{x}, \mathbf{x}') = f(\mathbf{x}) k(\mathbf{x}, \mathbf{x}') f(\mathbf{x}')$ is a kernel for any function $f$.
    \item \textbf{Difference is NOT generally a kernel}: counter-example by setting $\mathbf{x} = \mathbf{x}'$.
  \end{itemize}
  \item \textbf{Canonical kernels and their proofs}:
  \begin{itemize}
    \item \textbf{Linear}: $k(\mathbf{x}, \mathbf{x}') = \langle\mathbf{x}, \mathbf{x}'\rangle$. Trivially a kernel.
    \item \textbf{Polynomial}: $k(\mathbf{x}, \mathbf{x}') = (\langle\mathbf{x},\mathbf{x}'\rangle + c)^m$. Proof: binomial theorem expands to a conic sum of products of linear kernels.
    \item \textbf{Cosine}: $k(x, x') = \cos(x - x') = \cos(x)\cos(x') + \sin(x)\sin(x')$. Feature map $\phi(x) = (\cos x, \sin x)$. \emph{$\cos(x + x')$ is NOT a kernel} — counter-example at $x = x' = \pi/3$.
    \item \textbf{Exponential}: $k(\mathbf{x}, \mathbf{x}') = \exp(\langle\mathbf{x}, \mathbf{x}'\rangle)$. Proof: Taylor series of $e^z$ as a conic sum of polynomial kernels.
    \item \textbf{Gaussian / RBF}: $k(\mathbf{x}, \mathbf{x}') = \exp(-\|\mathbf{x}-\mathbf{x}'\|^2 / 2\sigma^2)$. Proof: rewrite using $\|\mathbf{x}-\mathbf{x}'\|^2 = \|\mathbf{x}\|^2 + \|\mathbf{x}'\|^2 - 2\langle\mathbf{x}, \mathbf{x}'\rangle$, then mapping rule.
  \end{itemize}
  \item \textbf{Kernel methods vs deep learning}: kernels embed inputs into a fixed feature space (chosen by the practitioner). Deep nets \emph{learn} the feature representation end-to-end. ``Kernels were the backbone of ML from early 90's to about a decade ago.''
\end{itemize}
\end{tldrbox}

% =====================================================================
\section{Where Wk29 fits}

Wk26 introduced kernels intuitively: the SVM dual depends only on $\phi(\mathbf{x}_i)^\top\phi(\mathbf{x}_j)$, so we can replace these inner products with a function $K(\mathbf{x}_i, \mathbf{x}_j)$ that computes the same value without ever forming $\phi$. Wk29 makes this rigorous and develops a construction calculus.

\textbf{Why the dual — the cost of the explicit dot product.} \examflag{}The dual objective contains $\sum_{i,j}\phi(\mathbf{x}_i)^\top\phi(\mathbf{x}_j)$ over every pair of training points. Forming $\phi$ explicitly means computing a dot product in the feature space for each pair — and a nonlinear feature transformation can blow the dimensionality up enormously (e.g. $10^3 \to 10^4 \to 10^5$ dimensions, possibly a million). The cost is therefore \emph{quadratic} in the number of training points $n$ and \emph{linear} in the (possibly exploding) feature dimension. Kernels exist precisely to compute that same inner product directly via $k(\mathbf{x}_i, \mathbf{x}_j)$, without ever forming $\phi$ — this is the core ``why'' of the topic.

\textbf{The $n$-vs-$d$ historical framing.} There are two scaling factors: the number of training samples $n$, and the data/feature dimensionality $d$. Decades ago $n$ was small (pre-big-data, no GPUs) — typically thousands, not millions — so the binding constraint was the high dimensionality of $\phi(\mathbf{x})$, which is exactly what kernels relieve.

\textbf{The core questions Wk29 answers:}
\begin{itemize}
  \item What \emph{is} a kernel formally? When does a feature space $\mathcal{H}$ and feature map $\phi$ exist?
  \item How do we build new kernels from existing ones?
  \item How do we \emph{prove} a given function is a kernel?
\end{itemize}

\textbf{Cross-references to Wk26} (recall, do not re-derive):
\begin{itemize}
  \item Soft-margin SVM, hinge loss, KKT conditions — Wk26 §2--§4.
  \item Dual SVM formulation, kernel SVM, prediction rule — Wk26 §6--§7.
  \item Polynomial and Gaussian/RBF kernels (introduced informally) — Wk26 §7.
\end{itemize}

% =====================================================================
\section{Formal kernel definition}

\subsection{The general definition}

\begin{defbox}[Kernel function]
For a non-empty set $\mathcal{X}$, a function $k : \mathcal{X}\times\mathcal{X} \to \mathbb{R}$ is a \textbf{kernel} if there exists a \emph{feature space} $\mathcal{H}$ (a Hilbert space) and a \emph{feature map} $\phi : \mathcal{X} \to \mathcal{H}$ such that
\[
k(\mathbf{x}, \mathbf{x}') \;=\; \langle\phi(\mathbf{x}), \phi(\mathbf{x}')\rangle_{\mathcal{H}} \qquad \forall \mathbf{x}, \mathbf{x}' \in \mathcal{X}.
\]
\end{defbox}

Three points to read carefully:
\begin{itemize}
  \item \textbf{$\mathcal{X}$ does not need to be a vector space.} It can be a set of documents, graphs, strings, images, time series — anything with a notion of ``individual instance.'' The kernel implicitly defines an inner-product structure even when no natural one exists in $\mathcal{X}$.
  \item \textbf{$\phi$ is the feature map}, $\mathcal{H}$ the feature space.
  \item \textbf{$\mathcal{H}$ must be a Hilbert space} — a vector space with a well-defined inner product. The technical conditions are not examinable; treat $\mathcal{H}$ as ``a vector space, possibly infinite-dimensional, in which inner products make sense.'' When $\mathcal{H}$ is infinite-dimensional the inner product is an infinite sum, and the Hilbert-space condition is precisely what guarantees that this sum converges to a finite number rather than diverging.
\end{itemize}

\begin{center}

\footnotesize\textit{A dataset that illustrates why a linear decision boundary is insufficient. \textbf{Blue crosses} and \textbf{red crosses} form two classes whose boundary is a winding, nonlinear contour --- no single straight line (or hyperplane in higher dimensions) can separate them. The \textbf{green-circled points} are the support vectors: the examples lying closest to, or straddling, the decision boundary and therefore determining its shape. This is the fundamental motivation for kernel methods --- by mapping the data into a richer feature space $\mathcal{H}$ via the feature map $\phi$, a linear hyperplane in $\mathcal{H}$ can implement the curved boundary visible here.}
\end{center}

\subsection{The inner-product axioms (non-examinable, for completeness)}

\begin{defbox}[Inner product --- the one examinable property]
The lecturer states explicitly that you will \emph{not} be asked to write down the full definition of an inner product. The only property to retain is \textbf{positive definiteness}: for an inner product $\langle\cdot, \cdot\rangle : V\times V \to \mathbb{R}$,
\[
\langle v, v\rangle \ge 0 \quad\text{for all } v\in V, \qquad \text{with equality iff } v = \mathbf{0}.
\]
\end{defbox}

\textbf{Examinable takeaway from this aside.} The first axiom ($\langle v, v\rangle \ge 0$) is the foundation for the ``$\mathbf{x} = \mathbf{x}'$ trick'' used to disprove proposed kernels.

\subsection{The non-uniqueness of feature maps}

For a given kernel, the feature map and Hilbert space are \textbf{not unique}. The simplest example: the linear kernel $k(x, x') = \langle x, x'\rangle$ on $\mathbb{R}$.

\begin{exbox}[Two valid feature maps for the linear kernel]
\textbf{Option 1.} $\phi_1(x) = x$, $\mathcal{H}_1 = \mathbb{R}$. Then $\langle\phi_1(x), \phi_1(x')\rangle = x\,x' = k(x, x')$. $\checkmark$\\[0.3em]
\textbf{Option 2.} $\phi_2(x) = (x/\sqrt 2,\, x/\sqrt 2)$, $\mathcal{H}_2 = \mathbb{R}^2$. Then
\[
\langle\phi_2(x), \phi_2(x')\rangle = (x/\sqrt 2)(x'/\sqrt 2) + (x/\sqrt 2)(x'/\sqrt 2) = x x' = k(x, x'). \;\checkmark
\]
\end{exbox}

Both feature maps yield the same kernel. \textbf{Two consequences:}
\begin{itemize}
  \item When asked to ``provide a feature map,'' any valid one earns marks.
  \item Two SVMs trained with the same kernel but conceptually different $\phi$ are computationally identical.
\end{itemize}

% =====================================================================
\section{Construction rules --- building kernels from kernels}

To prove a candidate function is a kernel, you can either (a) construct an explicit feature map (rare, hard for complex kernels), or (b) build the function from known kernels using a small set of construction rules. The construction-rule approach is what the lecture emphasises.

\subsection{Positive scalar multiple}

\begin{defbox}[Positive scalar multiple of a kernel is a kernel]
If $k$ is a kernel and $\alpha > 0$, then $\alpha k$ is also a kernel.
\end{defbox}

\begin{exbox}[\doflag{}Proof]
$k(\mathbf{x}, \mathbf{x}') = \langle\phi(\mathbf{x}), \phi(\mathbf{x}')\rangle_{\mathcal{H}}$ for some $\phi, \mathcal{H}$.
\[
\alpha k(\mathbf{x}, \mathbf{x}') = \alpha\langle\phi(\mathbf{x}), \phi(\mathbf{x}')\rangle_{\mathcal{H}} = \langle\sqrt{\alpha}\,\phi(\mathbf{x}),\,\sqrt{\alpha}\,\phi(\mathbf{x}')\rangle_{\mathcal{H}}.
\]
The new feature map is $\sqrt{\alpha}\,\phi$ in the same Hilbert space.\hfill\textbf{QED}
\end{exbox}

\textbf{Why $\alpha > 0$?} If $\alpha < 0$ then $\alpha\langle\phi(\mathbf{x}), \phi(\mathbf{x})\rangle_{\mathcal{H}} \le 0$ for $\phi(\mathbf{x}) \ne \mathbf{0}$, violating positive-definiteness. \textbf{Negative-scalar multiples are not kernels.}

\subsection{Conic sum (positive linear combination)}

\begin{defbox}[Conic sum of kernels is a kernel]
If $k_1, \ldots, k_K$ are kernels and $\alpha_1, \ldots, \alpha_K > 0$, then $\sum_{j=1}^K \alpha_j k_j$ is a kernel.
\end{defbox}

\begin{exbox}[\doflag{}Proof]
Each $k_j(\mathbf{x}, \mathbf{x}') = \langle\phi_j(\mathbf{x}), \phi_j(\mathbf{x}')\rangle_{\mathcal{H}_j}$. Then
\[
\sum_j \alpha_j k_j(\mathbf{x}, \mathbf{x}') = \sum_j \langle\sqrt{\alpha_j}\phi_j(\mathbf{x}),\,\sqrt{\alpha_j}\phi_j(\mathbf{x}')\rangle_{\mathcal{H}_j}.
\]
Define the augmented feature map $\hat\phi(\mathbf{x}) = (\sqrt{\alpha_1}\phi_1(\mathbf{x}), \ldots, \sqrt{\alpha_K}\phi_K(\mathbf{x}))$ in the direct sum $\hat{\mathcal{H}} = \mathcal{H}_1 \oplus \cdots \oplus \mathcal{H}_K$. Then
\[
\sum_j \alpha_j k_j(\mathbf{x}, \mathbf{x}') = \langle\hat\phi(\mathbf{x}), \hat\phi(\mathbf{x}')\rangle_{\hat{\mathcal{H}}}.
\]
The conic sum is therefore an inner product in the concatenated feature space.\hfill\textbf{QED}
\end{exbox}

\textbf{Concatenation as direct sum.} The new feature space is a stacked vector — features of $k_1$ followed by features of $k_2$, etc. ``$\mathbb{R}^2 \oplus \mathbb{R}^1 = \mathbb{R}^3$.''

\subsection{Product (without proof)}

\begin{defbox}[Product of kernels is a kernel]
If $k_1, k_2$ are kernels, then $k_1(\mathbf{x}, \mathbf{x}')\cdot k_2(\mathbf{x}, \mathbf{x}')$ is a kernel.
\end{defbox}

The lecture states this fact without proof (the proof requires tensor-product Hilbert spaces). Take it as given for construction-style problems.

\subsection{Mapping wrap}

\begin{defbox}[Function-times-kernel-times-function is a kernel]
For any function $f : \mathcal{X} \to \mathbb{R}$ and any kernel $k$,
\[
\hat k(\mathbf{x}, \mathbf{x}') = f(\mathbf{x})\,k(\mathbf{x}, \mathbf{x}')\,f(\mathbf{x}')
\]
is also a kernel.
\end{defbox}

\begin{exbox}[\doflag{}Proof]
$k(\mathbf{x}, \mathbf{x}') = \langle\phi(\mathbf{x}), \phi(\mathbf{x}')\rangle_{\mathcal{H}}$.
\[
f(\mathbf{x})k(\mathbf{x}, \mathbf{x}')f(\mathbf{x}') = f(\mathbf{x})\langle\phi(\mathbf{x}), \phi(\mathbf{x}')\rangle_{\mathcal{H}} f(\mathbf{x}') = \langle f(\mathbf{x})\phi(\mathbf{x}),\,f(\mathbf{x}')\phi(\mathbf{x}')\rangle_{\mathcal{H}}.
\]
The new feature map is $f(\mathbf{x})\phi(\mathbf{x})$ in the \emph{same} Hilbert space $\mathcal{H}$.\hfill\textbf{QED}
\end{exbox}

\textbf{Use case.} Whenever a candidate kernel can be factored as $f(\mathbf{x}) g(\langle\mathbf{x}, \mathbf{x}'\rangle, \ldots) f(\mathbf{x}')$ with $g$ already a kernel, the candidate is too. The Gaussian kernel proof below uses this trick.

\subsection{Difference is not generally a kernel}

\begin{defbox}[$k_1 - k_2$ is generally NOT a kernel]
For two kernels $k_1, k_2$, the difference $k_1 - k_2$ is not necessarily a kernel.
\end{defbox}

\begin{exbox}[Counterexample sketch]
Suppose for some $\mathbf{x}\in\mathcal{X}$, $k_1(\mathbf{x}, \mathbf{x}) - k_2(\mathbf{x}, \mathbf{x}) < 0$. (If not, swap the roles to get the inequality the other way.)\\[0.3em]
If $k_1 - k_2$ were a kernel, there would exist $\phi$ and $\mathcal{H}$ such that
\[
k_1(\mathbf{x}, \mathbf{x}') - k_2(\mathbf{x}, \mathbf{x}') = \langle\phi(\mathbf{x}), \phi(\mathbf{x}')\rangle_{\mathcal{H}}.
\]
At $\mathbf{x} = \mathbf{x}'$, the right-hand side is $\|\phi(\mathbf{x})\|^2_{\mathcal{H}} \ge 0$ (positive-definiteness). But the left-hand side is negative — contradiction. So $k_1 - k_2$ is not a kernel for that pair.\hfill\textbf{QED}
\end{exbox}

\textbf{Memorise this trick.} The ``$\mathbf{x} = \mathbf{x}'$ argument'' is the canonical way to show a candidate is \emph{not} a kernel. We use it again to disprove $\cos(x + x')$.

% =====================================================================
\section{Canonical kernels and their proofs}

\subsection{The polynomial kernel}

\begin{defbox}[Polynomial kernel]
For $\mathbf{x}, \mathbf{x}' \in \mathbb{R}^d$, $c \ge 0$, $m \ge 1$ integer:
\[
k(\mathbf{x}, \mathbf{x}') = (\langle\mathbf{x}, \mathbf{x}'\rangle + c)^m.
\]
With $m = 1, c = 0$: \textbf{linear kernel} $k(\mathbf{x}, \mathbf{x}') = \langle\mathbf{x}, \mathbf{x}'\rangle$.
\end{defbox}

\begin{exbox}[\doflag{}Proof that the polynomial kernel is a kernel]
Apply the binomial theorem:
\[
(\langle\mathbf{x}, \mathbf{x}'\rangle + c)^m = \sum_{i=0}^m \binom{m}{i} c^{m-i} \langle\mathbf{x}, \mathbf{x}'\rangle^i.
\]
Each term consists of:
\begin{itemize}
  \item $\binom{m}{i} c^{m-i} \ge 0$ (positive coefficient).
  \item $\langle\mathbf{x}, \mathbf{x}'\rangle^i$, a product of $i$ copies of the linear kernel $\langle\mathbf{x}, \mathbf{x}'\rangle$. By the product-of-kernels rule, each $\langle\mathbf{x}, \mathbf{x}'\rangle^i$ is a kernel.
\end{itemize}
The whole sum is a conic sum of kernels with positive coefficients, hence a kernel.\hfill\textbf{QED}
\end{exbox}

\begin{center}

\footnotesize\textit{A concrete demonstration of the feature-map idea. \textbf{Left}: the original 2-D input space spanned by $x^{(1)}$ and $x^{(2)}$, where the \textbf{red} and \textbf{blue} clusters interleave and no straight line separates them. \textbf{Right}: the same data after introducing a third feature $x^{(1)} x^{(2)}$ (the product of the two original features), lifting each point into 3-D. In this richer space the two classes lie on opposite sides of the grey separating plane --- the data has become linearly separable. This is exactly what the polynomial kernel $k(\mathbf{x}, \mathbf{x}') = (\langle\mathbf{x}, \mathbf{x}'\rangle + c)^m$ achieves implicitly: the binomial expansion corresponds to the set of all monomial features up to degree $m$, so the dual SVM never needs to form $\phi(\mathbf{x})$ explicitly --- the kernel function computes the necessary inner products directly.}
\end{center}

\textbf{What the polynomial feature space concretely is.} For a 2-D input $\mathbf{x} = (x^{(1)}, x^{(2)})$, the degree-2 polynomial kernel corresponds to the explicit feature map
\[
\phi(\mathbf{x}) = \big(1,\; x^{(1)},\; x^{(2)},\; x^{(1)}x^{(2)},\; (x^{(1)})^2,\; (x^{(2)})^2\big).
\]
A degree-$m$ polynomial kernel gives all $m$-th-order \textbf{feature interactions}: for $m = 2$ these are the pairwise products and squares above; for larger $m$ the feature space contains correspondingly higher-order interactions among the input coordinates.

\subsection{The cosine kernel --- and what's not a kernel}

\begin{defbox}[Cosine kernel]
For $x, x' \in \mathbb{R}$:
\[
k_{\text{cos}}(x, x') = \cos(x - x').
\]
\end{defbox}

\begin{exbox}[\doflag{}Proof: cosine kernel via the cosine subtraction formula]
Use $\cos(x - x') = \cos x\cos x' + \sin x\sin x'$:
\[
k_{\text{cos}}(x, x') = \cos x\cos x' + \sin x\sin x' = \langle\phi(x), \phi(x')\rangle
\]
with $\phi(x) = (\cos x, \sin x) \in \mathbb{R}^2$.\hfill\textbf{QED}
\end{exbox}

\textbf{Note the explicit feature map.} The cosine kernel maps the real line to the unit circle in $\mathbb{R}^2$.

\textbf{Compare: $\cos(x + x')$ is NOT a kernel.}

\begin{exbox}[\doflag{}Counterexample showing $\cos(x + x')$ is not a kernel]
Set $x = x' = \pi/3$:
\[
\cos(x + x') = \cos(2\pi/3) = -0.5 < 0.
\]
But $\langle\phi(x), \phi(x)\rangle = \|\phi(x)\|^2 \ge 0$ for any feature map. Contradiction. So no feature map exists.\hfill\textbf{QED}
\end{exbox}

\textbf{The lesson.} Small algebraic differences (sum vs difference inside $\cos$) can flip a kernel into a non-kernel. The $\mathbf{x} = \mathbf{x}'$ trick disposes of candidates quickly.

\subsection{The exponential kernel}

\begin{defbox}[Exponential kernel]
For $\mathbf{x}, \mathbf{x}' \in \mathbb{R}^d$:
\[
k_{\text{exp}}(\mathbf{x}, \mathbf{x}') = \exp(\langle\mathbf{x}, \mathbf{x}'\rangle).
\]
\end{defbox}

\begin{exbox}[\doflag{}Proof: exponential kernel via Taylor series]
Use $e^z = \sum_{i=0}^\infty z^i / i!$:
\[
\exp(\langle\mathbf{x}, \mathbf{x}'\rangle) = \sum_{i=0}^\infty \frac{\langle\mathbf{x}, \mathbf{x}'\rangle^i}{i!}.
\]
Each term:
\begin{itemize}
  \item $1/i! \ge 0$ (positive coefficient).
  \item $\langle\mathbf{x}, \mathbf{x}'\rangle^i$ is a kernel (product of linear kernels).
\end{itemize}
The infinite sum is a conic sum of kernels — a kernel.\hfill\textbf{QED}
\end{exbox}

\textbf{Why this matters.} The exponential kernel feature space is the infinite-dimensional space of all polynomials in $\mathbf{x}$. The Taylor expansion makes this concrete.

\subsection{The Gaussian / RBF kernel}

\begin{defbox}[Gaussian / RBF kernel]
For $\mathbf{x}, \mathbf{x}' \in \mathbb{R}^d$ and fixed bandwidth $\sigma > 0$:
\[
k_{\text{Gauss}}(\mathbf{x}, \mathbf{x}') = \exp\!\left(-\frac{\|\mathbf{x} - \mathbf{x}'\|^2}{2\sigma^2}\right).
\]
\textbf{Translation invariant} — depends only on $\mathbf{x} - \mathbf{x}'$.
\end{defbox}

\begin{exbox}[\doflag{}Proof: Gaussian kernel]
Expand $\|\mathbf{x} - \mathbf{x}'\|^2 = \|\mathbf{x}\|^2 + \|\mathbf{x}'\|^2 - 2\langle\mathbf{x}, \mathbf{x}'\rangle$:
\[
\exp\!\left(-\frac{\|\mathbf{x} - \mathbf{x}'\|^2}{2\sigma^2}\right) = \exp\!\left(-\frac{\|\mathbf{x}\|^2}{2\sigma^2}\right)\cdot \exp\!\left(\frac{\langle\mathbf{x}, \mathbf{x}'\rangle}{\sigma^2}\right) \cdot \exp\!\left(-\frac{\|\mathbf{x}'\|^2}{2\sigma^2}\right).
\]
\textbf{Read this as a mapping wrap.} Define $f(\mathbf{x}) = \exp(-\|\mathbf{x}\|^2 / 2\sigma^2)$. The middle factor is the exponential kernel $\exp(\langle\mathbf{x}, \mathbf{x}'\rangle / \sigma^2)$ — itself a kernel by the previous proof (apply the $\sigma^{-2}$ rescaling first). So
\[
k_{\text{Gauss}}(\mathbf{x}, \mathbf{x}') = f(\mathbf{x})\cdot k_{\text{exp}}(\mathbf{x}, \mathbf{x}'/\sigma^2)\cdot f(\mathbf{x}').
\]
By the mapping-wrap rule, this is a kernel.\hfill\textbf{QED}
\end{exbox}

\textbf{The infinite feature space.} The exponential kernel's feature space is infinite-dimensional. The mapping wrap by $f$ doesn't change that. \textbf{The Gaussian kernel's feature space is therefore infinite-dimensional.} You cannot ever form $\phi(\mathbf{x})$ explicitly — but the kernel function gives you the inner product directly. This is why kernels are so powerful.

\subsection{\texorpdfstring{The bandwidth $\sigma$ as a knob}{The bandwidth sigma as a knob}}

\begin{tabularx}{\linewidth}{@{}lXX@{}}
\toprule
\thead{Bandwidth $\sigma$} & \thead{Behaviour} & \thead{SVM consequence} \\
\midrule
Small $\sigma$ & each training point's influence is highly localised & overfitting risk; complex, wiggly decision boundaries \\
Large $\sigma$ & each training point's influence is broad        & under-fitting; smooth decision boundaries, more like linear \\
\bottomrule
\end{tabularx}

\textbf{Cross-validate $\sigma$.} The bandwidth is a hyperparameter to tune, like $C$ in soft-margin SVM. ``RBF kernel with $\gamma = 0.1, C = 1$'' (where $\gamma = 1/(2\sigma^2)$) is the typical sklearn default tune-from.

% =====================================================================
\section{Comparison with deep learning --- where do the features come from?}

\textbf{In an SVM with a kernel}, the feature map $\phi$ is fixed. You choose it (explicitly via $\phi$, or implicitly via the kernel) before seeing the data. The model learns a linear separator in the feature space, but the feature space is given.

\textbf{In a deep network}, the feature map $\phi$ is learned end-to-end. Each layer produces a new representation; the optimisation finds a feature space in which the data is linearly separable (or at least cleanly classifiable by the final layer).

\begin{tabularx}{\linewidth}{@{}lXX@{}}
\toprule
& \thead{Kernel SVM} & \thead{Deep neural network} \\
\midrule
Feature representation & fixed by kernel choice                  & learned from data \\
Optimisation problem  & convex (QP)                               & non-convex (many local optima) \\
Sample efficiency     & high — works on small datasets          & lower — typically needs many samples \\
Computational scaling & $O(n^2)$ or worse in $n$ training points & $O(n)$ per pass; many passes \\
Source of non-linearity & non-linear kernel feature map         & non-linear activation functions \\
Era of dominance      & 1990s --- mid-2010s                       & mid-2010s --- present \\
\bottomrule
\end{tabularx}

\textbf{Lecturer's framing.} ``In the context of SVM, the role of non-linear features (arising out of product or other similar operations) explicitly gives non-linearity. The same effect can be achieved implicitly in a computationally efficient manner by using kernels. However, the feature transformation corresponding to the above operations is fixed and is relatively hard to learn in problem-dependent manner. Deep learning is about learning feature representations as part of a non-convex problem.''

\textbf{Both have their roles.} Kernel SVMs remain the dominant choice on small structured datasets (medical, financial, scientific). Deep nets dominate when massive labelled data is available (vision, language).

% =====================================================================
\section{Exam-mode answers}

\begin{exambox}[{\examflag{}(MCQ Past paper Q2(6)): A major advantage of kernel methods is that they can handle: (a) high-dim data, (b) non-linear relationships, (c) neither, (d) both.}]
\textit{\textbf{Correct: (d) Both.} Kernel methods enable both simultaneously. (a) Kernels operate in implicit feature spaces that may be high-dimensional or even infinite-dimensional (RBF) without ever explicitly forming the features. (b) The non-linear feature map $\phi$ (or equivalently the non-linear kernel function) produces non-linear decision boundaries in the original input space. (2 marks)}
\end{exambox}

\begin{exambox}[{Concept: State the formal definition of a kernel.}]
\textit{$k : \mathcal{X}\times\mathcal{X}\to\mathbb{R}$ is a \textbf{kernel} if there exists a Hilbert space $\mathcal{H}$ (the feature space) and a function $\phi : \mathcal{X}\to\mathcal{H}$ (the feature map) such that for all $\mathbf{x}, \mathbf{x}' \in \mathcal{X}$:}
\[
k(\mathbf{x}, \mathbf{x}') = \langle\phi(\mathbf{x}), \phi(\mathbf{x}')\rangle_{\mathcal{H}}.
\]
\textit{$\mathcal{X}$ does not need to be a vector space — it can be any non-empty set (documents, graphs, strings).}
\end{exambox}

\begin{exambox}[{\doflag{}Concept: Show that the polynomial kernel $k(\mathbf{x}, \mathbf{x}') = (\langle\mathbf{x}, \mathbf{x}'\rangle + c)^m$ is a kernel.}]
\textit{Apply the binomial theorem:}
\[
(\langle\mathbf{x}, \mathbf{x}'\rangle + c)^m = \sum_{i=0}^m \binom{m}{i} c^{m-i} \langle\mathbf{x}, \mathbf{x}'\rangle^i.
\]
\textit{The coefficient $\binom{m}{i}c^{m-i} \ge 0$ for $c \ge 0$. The factor $\langle\mathbf{x}, \mathbf{x}'\rangle^i$ is a product of $i$ copies of the linear kernel and is therefore a kernel (product of kernels rule). The whole expression is a conic sum of kernels with positive coefficients, hence a kernel.}
\end{exambox}

\begin{exambox}[{\doflag{}Concept: Show that the Gaussian kernel is a kernel.}]
\textit{$k_{\text{Gauss}}(\mathbf{x}, \mathbf{x}') = \exp(-\|\mathbf{x}-\mathbf{x}'\|^2/2\sigma^2)$. Use $\|\mathbf{x}-\mathbf{x}'\|^2 = \|\mathbf{x}\|^2 + \|\mathbf{x}'\|^2 - 2\langle\mathbf{x}, \mathbf{x}'\rangle$:}
\[
k_{\text{Gauss}}(\mathbf{x}, \mathbf{x}') = \underbrace{\exp(-\|\mathbf{x}\|^2/2\sigma^2)}_{f(\mathbf{x})}\cdot \underbrace{\exp(\langle\mathbf{x}, \mathbf{x}'\rangle/\sigma^2)}_{\text{exponential kernel}}\cdot \underbrace{\exp(-\|\mathbf{x}'\|^2/2\sigma^2)}_{f(\mathbf{x}')}.
\]
\textit{The middle factor is the exponential kernel (a kernel by Taylor expansion of $e^z$ as a conic sum of polynomial kernels). The outer factors are the same function $f$ evaluated at $\mathbf{x}$ and $\mathbf{x}'$. By the mapping-wrap rule ($f(\mathbf{x})k(\mathbf{x},\mathbf{x}')f(\mathbf{x}')$ is a kernel), the whole expression is a kernel.}
\end{exambox}

\begin{exambox}[{\doflag{}Concept: Is $\cos(x + x')$ a kernel? Justify.}]
\textit{\textbf{No.} Set $x = x' = \pi/3$: $\cos(x + x') = \cos(2\pi/3) = -0.5 < 0$. If $\cos(x + x')$ were a kernel, $\cos(x + x)\big|_{x = x'} = \langle\phi(x), \phi(x)\rangle = \|\phi(x)\|^2 \ge 0$. But $-0.5 < 0$, contradicting positive-definiteness of the inner product. Hence $\cos(x + x')$ is not a kernel.}
\end{exambox}

\begin{exambox}[{\doflag{}Concept: Show that $\cos(x - x')$ \emph{is} a kernel and provide an explicit feature map.}]
\textit{Use $\cos(x - x') = \cos x \cos x' + \sin x \sin x'$. Setting $\phi(x) = (\cos x, \sin x)$ in $\mathcal{H} = \mathbb{R}^2$:}
\[
\langle\phi(x), \phi(x')\rangle = \cos x\cos x' + \sin x\sin x' = \cos(x - x').
\]
\textit{So $\cos(x - x')$ is a kernel with feature map $\phi(x) = (\cos x, \sin x)$.}
\end{exambox}

\begin{exambox}[{Concept: What construction rules can you use to combine kernels?}]
\textit{\textbf{(1)} Positive scalar multiple: $\alpha k$ for $\alpha > 0$.}\\
\textit{\textbf{(2)} Conic sum: $\sum_j \alpha_j k_j$ for $\alpha_j > 0$.}\\
\textit{\textbf{(3)} Product: $k_1 \cdot k_2$.}\\
\textit{\textbf{(4)} Mapping wrap: $f(\mathbf{x})\,k(\mathbf{x}, \mathbf{x}')\,f(\mathbf{x}')$ for any function $f : \mathcal{X}\to\mathbb{R}$.}\\[0.3em]
\textit{\textbf{Difference is NOT generally a kernel.} Counter-example by the $\mathbf{x} = \mathbf{x}'$ trick: $k_1(\mathbf{x}, \mathbf{x}) - k_2(\mathbf{x}, \mathbf{x})$ can be negative, contradicting $\langle\phi(\mathbf{x}), \phi(\mathbf{x})\rangle \ge 0$.}
\end{exambox}

\begin{exambox}[{Concept: Compare kernel SVMs and deep neural networks on feature representation.}]
\textit{\textbf{Kernel SVM.} Feature map $\phi$ is fixed in advance — chosen by the user (polynomial degree, RBF bandwidth, etc.). The model learns a linear separator in this fixed feature space. The optimisation is convex.}\\[0.3em]
\textit{\textbf{Deep neural network.} The feature representation is \emph{learned} end-to-end as part of training. Each layer produces a new feature space; the final layer separates linearly in the topmost representation. The optimisation is non-convex.}\\[0.3em]
\textit{Trade-off: kernels are sample-efficient but rely on a good fixed kernel choice; deep nets need more data but can learn problem-specific representations. Kernels dominated ML from the early 90s to the mid-2010s; deep learning since.}
\end{exambox}

% =====================================================================
\section*{Key equations cheat sheet}
\addcontentsline{toc}{section}{Key equations cheat sheet}

\begin{tabularx}{\linewidth}{@{}lX@{}}
\toprule
\thead{Object} & \thead{Formula / fact} \\
\midrule
Kernel definition         & $k(\mathbf{x}, \mathbf{x}') = \langle\phi(\mathbf{x}), \phi(\mathbf{x}')\rangle_{\mathcal{H}}$ for some $\phi, \mathcal{H}$ \\
Linear kernel             & $k(\mathbf{x}, \mathbf{x}') = \langle\mathbf{x}, \mathbf{x}'\rangle$ \\
Polynomial kernel         & $k(\mathbf{x}, \mathbf{x}') = (\langle\mathbf{x}, \mathbf{x}'\rangle + c)^m$, $c \ge 0$, $m \ge 1$ \\
Cosine kernel             & $k(x, x') = \cos(x - x')$, $\phi(x) = (\cos x, \sin x)$ \\
NOT a kernel              & $\cos(x + x')$ — counter-example at $x = \pi/3$ \\
Exponential kernel        & $k(\mathbf{x}, \mathbf{x}') = \exp(\langle\mathbf{x}, \mathbf{x}'\rangle)$ \\
Gaussian / RBF kernel     & $k(\mathbf{x}, \mathbf{x}') = \exp(-\|\mathbf{x}-\mathbf{x}'\|^2 / 2\sigma^2)$ \\
Positive scalar           & $\alpha k$ kernel for $\alpha > 0$ \\
Conic sum                 & $\sum_j \alpha_j k_j$ kernel for $\alpha_j > 0$ \\
Product                   & $k_1\cdot k_2$ kernel \\
Mapping wrap              & $f(\mathbf{x})k(\mathbf{x}, \mathbf{x}')f(\mathbf{x}')$ kernel \\
Difference                & $k_1 - k_2$ \emph{not generally} a kernel \\
Disproof technique        & set $\mathbf{x} = \mathbf{x}'$ and check $k(\mathbf{x}, \mathbf{x}) \ge 0$ \\
\bottomrule
\end{tabularx}

% =====================================================================
\section*{Pitfalls and traps consolidated}
\addcontentsline{toc}{section}{Pitfalls and traps consolidated}

\begin{itemize}
  \item \textbf{$\mathcal{X}$ is not necessarily a vector space.} Kernels work on documents, graphs, strings, time series. Don't assume Euclidean.
  \item \textbf{Feature maps are non-unique.} Two different $\phi$ can give the same kernel. ``Provide a feature map'' has multiple valid answers.
  \item \textbf{Negative scalar multiples are not kernels.} The proof requires $\sqrt{\alpha}$, which is real only for $\alpha \ge 0$.
  \item \textbf{Difference is not a kernel.} The $\mathbf{x} = \mathbf{x}'$ trick disposes of this quickly.
  \item \trapflag{}\textbf{$\cos(x + x')$ vs $\cos(x - x')$.} The first is not a kernel (sign flip at $x = \pi/3$); the second is (cosine subtraction formula). Easy to get backwards.
  \item \textbf{Gaussian kernel's feature space is infinite-dimensional.} You cannot compute $\phi(\mathbf{x})$ explicitly — only the kernel function. This is the magic.
  \item \textbf{Bandwidth $\sigma$ is a hyperparameter.} Must be tuned. Too small $\Rightarrow$ overfit; too large $\Rightarrow$ underfit. Cross-validate.
  \item \textbf{Past paper Q2(6): both (a) and (b).} Don't be tempted to pick only one.
  \item \textbf{Construction rules require the right sign.} Conic sum: $\alpha_j > 0$. Mapping wrap: $f$ can be any function (signs don't matter because $f(\mathbf{x})f(\mathbf{x}')$ pairs up symmetrically).
  \item \textbf{Prove a candidate is a kernel} by either constructing $\phi$ explicitly or building from known kernels via construction rules. The latter is usually cleaner.
\end{itemize}

% =====================================================================
\section*{Self-check questions}
\addcontentsline{toc}{section}{Self-check questions}

\begin{enumerate}
  \item State the formal definition of a kernel. What conditions must $\mathcal{H}$ satisfy?
  \item Give an example of a non-vector-space $\mathcal{X}$ on which kernels make sense.
  \item Provide two distinct valid feature maps for the linear kernel $k(x, x') = \langle x, x'\rangle$.
  \item State and prove the positive-scalar-multiple rule for kernels.
  \item State and prove the conic-sum rule for kernels.
  \item State the mapping-wrap rule (function $\times$ kernel $\times$ function) and prove it.
  \item Show that the difference $k_1 - k_2$ is not generally a kernel. Use the $\mathbf{x} = \mathbf{x}'$ trick.
  \item Prove that the polynomial kernel $k(\mathbf{x}, \mathbf{x}') = (\langle\mathbf{x}, \mathbf{x}'\rangle + c)^m$ is a kernel using the binomial theorem.
  \item Prove that $\cos(x - x')$ is a kernel and provide an explicit feature map.
  \item Prove that $\cos(x + x')$ is \emph{not} a kernel by constructing a counter-example.
  \item Prove that the exponential kernel $\exp(\langle\mathbf{x}, \mathbf{x}'\rangle)$ is a kernel using the Taylor series of $e^z$.
  \item Prove that the Gaussian / RBF kernel is a kernel using the decomposition $\|\mathbf{x}-\mathbf{x}'\|^2 = \|\mathbf{x}\|^2 + \|\mathbf{x}'\|^2 - 2\langle\mathbf{x}, \mathbf{x}'\rangle$ and the mapping-wrap rule.
  \item Past paper Q2(6): kernel methods can handle (a) high-dim data, (b) non-linear relationships. Which is correct? Justify.
  \item Compare kernel SVMs to deep neural networks on (a) feature representation source, (b) optimisation type, (c) sample efficiency.
  \item Why is the Gaussian kernel's bandwidth $\sigma$ a hyperparameter requiring tuning? What are the consequences of $\sigma$ too small or too large?
\end{enumerate}

\end{document}
